\chapter{Manufacturing Validation Reference}
\label{app:manufacturing-reference}

This appendix is a lookup reference for program-gate manufacturing evidence,
laser-bearing production-control checklists, gauge and station templates, and
supplier first-article packages. It does not build the manufacturing-validation
argument. That narrative lives in \Cref{ch:manufacturing}: freeze the
production reference, build representative populations, preserve genealogy,
validate the measurement system, map yield, establish capability and
guardbands, choose controls, run SPC and reaction plans, manage changes, and
feed escapes back into controls.

\textit{Read first:} which gate asks which manufacturing question; how control
classes differ.

\textit{Reference:} ATP checklist; FAIR and supplier evidence; Rapid Interview
Checks.

\section{Program-gate manufacturing evidence}
\label{sec:mfg-npi-reference}

Program-phase labels (EVT, DVT, PVT, MP) are owned in
\Cref{sec:phase-labels,tab:phase-readiness-map}. They are not a second
lifecycle. The table below is a manufacturing-evidence lookup: which release
call the gate evidence supports when the program uses those names.

\begin{table*}[ht]
\refstepcounter{table}\label{tab:npi}%
\centering
\setlength{\tabcolsep}{4pt}%
\footnotesize
\renewcommand{\arraystretch}{1.25}%
\begin{tabular}{@{}p{0.10\linewidth}p{0.22\linewidth}p{0.36\linewidth}p{0.26\linewidth}@{}}
\toprule
Gate & Question & Representative evidence & Release decision \\
\midrule
EVT & Does it operate at all? & First light; CMIS bring-up; basic LIV/SMSR/RIN; one link closes BER & Continue / redesign integration \\
\addlinespace[1.25em]
DVT & Does it meet spec across corners? & Full ATP at $T$/$V$; prod-rep corners; stress plan + FIT model frozen & Enter qual / PVT / hold \\
\addlinespace[1.25em]
Qual & Env / reliability evidence ready? & Named mechanisms; sample plan; confidence (\Cref{sec:tree-qual-evidence}) & Enter PVT / hold \\
\addlinespace[1.25em]
PVT & Is it buildable at yield? & Multi-lot first-pass yield; MSA; ATP coverage; process capability; traceability; supplier readiness; validated rework; FAIR; production-host bring-up & Enter pilot / hold \\
\addlinespace[1.25em]
Pilot & Do assumptions hold in a bounded field trial? & Known serials/lots; enhanced telemetry; exit/rollback criteria & Open MP / restrict \\
\addlinespace[1.25em]
MP & Is quality sustained? & Steady DPPM; owned RMA Pareto; ECO control; fleet feedback & Keep shipping / CAPA / restrict \\
\bottomrule
\end{tabular}
\par\medskip
{\raggedright\footnotesize\textbf{Table~\thetable.} NPI gates as a manufacturing
evidence lookup. Pilot sits between PVT and MP. MP is sustained control, not
fleet monitoring alone. Phase definitions:
\Cref{sec:phase-labels}. Argument structure: \Cref{ch:manufacturing}.\par}
\end{table*}
\FloatBarrier

\section{Laser-bearing production-control checklist}
\label{sec:atp-laser-read}

Every second of ATP costs line capacity. Every omitted screen leaves risk. Buy
evidence in categories that cannot substitute for each other.
\Cref{tab:atp-laser} is a working checklist for an EML pluggable or an ELSFP CW
module. Customize limits from the datasheet and the link budget; do not invent
numbers in the ATP itself. Control-selection logic lives in
\Cref{sec:hvm-test,tab:mfg-control-types}.

\paragraph{Every-unit fast screens.}
Power class, CMIS state and bring-up, basic lane operation, and selected
electrical or optical checks that are cheap enough for 100\% coverage. These
catch dead units and firmware fails. They do not establish life.

\paragraph{Sampled expensive measurements.}
TDECQ, intrinsic and stressed RIN, detailed spectrum/SMSR, thermal corners, and
longer BER dwell. Use lot sample or audit when the signature is expensive and
correlated to escape risk.

\paragraph{Qualification-only stresses.}
HTOL, humidity, extensive cycling, and destructive analysis gather life or
mechanism evidence. They are not per-unit ship screens
(\Cref{sec:gr468,sec:tree-qual-evidence}).

\paragraph{Process controls.}
SPC on LIV, SMSR, RIN, TDECQ, and mate-cycle yield; golden-unit tracking; gauge
R\&R; incoming quality; first-article / FAIR after tooling or site change
(\Cref{sec:gauge-rr}).

\paragraph{Fleet controls.}
Telemetry, lot traceability, RMA codes split by mechanism, and recurrence
monitoring (\Cref{sec:fleet-triage}). These catch what ATP never saw.

A fast production power check does not establish life. HTOL does not prove every
shipped unit has correct firmware. SPC does not detect an unmeasured mechanism.
Source-level LIV, SMSR, and wavelength measurements can detect selected material
or device shifts when those measurements are directly available and correlated
to product risk. In a closed module, use a validated module-level proxy,
supplier evidence, sampled audit, or genealogy-based control; do not claim
internal measurement coverage when the production architecture does not expose
it. RIN and ORL protect the reflection environment; EAM/DCA protects Tx quality
on EML paths; CMIS protects field evidence. Burn-in may screen a demonstrated
infant-mortality population. ESD robustness is primarily established through
design qualification and handling controls; a finished-unit production screen
may not reliably detect latent ESD damage. Thermal class protects the derate
claim.

\begin{table*}[ht]
\refstepcounter{table}\label{tab:atp-laser}%
\centering
\setlength{\tabcolsep}{2.5pt}%
\footnotesize
\renewcommand{\arraystretch}{1.2}%
\begin{tabular}{@{}p{0.18\linewidth}p{0.16\linewidth}p{0.14\linewidth}p{0.22\linewidth}p{0.22\linewidth}@{}}
\toprule
Item & Method & Control class & Pass intent & Ties to \\
\midrule
LIV ($I_\mathrm{th}$, slope, kink) & SMU + power meter & 100\% ATP (source) / module proxy & kink-free bias window & wear-out, derate \\
\addlinespace[1.1em]
SMSR & OSA & 100\% ATP (source) / lot sample & single-mode vs.\ floor & modal noise \\
\addlinespace[1.1em]
Intrinsic RIN & PD + ESA & Lot sample / FA & quiet source floor & BER floor budget \\
\addlinespace[1.1em]
Stressed $\mathrm{RIN}_x\mathrm{OMA}$ & PD + ESA @ ORL & Lot sample / 100\% if escape & stressed Tx metric & named PMD \\
\addlinespace[1.1em]
Wavelength / grid & OSA / wavemeter & 100\% ATP or sample & channel ID & WDM lock \\
\addlinespace[1.1em]
Optical power class & power meter & 100\% ATP & class met & link budget \\
\addlinespace[1.1em]
EAM / TDECQ (EML) & bias + DCA & Lot sample / audit & ER, RLM, TDECQ & Tx quality \\
\addlinespace[1.1em]
CMIS / TWI bring-up & host / CMIS tool & 100\% ATP & state machine & telemetry \\
\addlinespace[1.1em]
Connector / ORL & mate + ORL meter & Periodic audit / sample & cycles + endface & packaging \\
\addlinespace[1.1em]
Burn-in (infant) & production screen & 100\% or lot sample & infant culled & not HTOL life \\
\addlinespace[1.1em]
HTOL life evidence & accelerated life & Qualification only & mechanism + FIT claim & GR-468 \\
\addlinespace[1.1em]
Driver/TIA ESD & JESD47 report & Qualification / audit & rating on file & IC reliability \\
\addlinespace[1.1em]
Thermal class & chamber & Lot sample / SPC & LIV/RIN/CMIS pass & derate \\
\addlinespace[1.1em]
Process monitors & SPC charts & SPC / process & drift detection & escape prevention \\
\addlinespace[1.1em]
Fleet cohort metrics & telemetry & Fleet telemetry & trend / alarm & recurrence \\
\bottomrule
\end{tabular}
\par\medskip
{\raggedright\footnotesize\textbf{Table~\thetable.} Control checklist for
laser-bearing modules (EML or ELSFP). Control class separates 100\% ATP, lot
sample, audit, qualification-only, SPC, and fleet telemetry. Intrinsic RIN and
stressed $\mathrm{RIN}_x\mathrm{OMA}$ are different metrics.\par}
\end{table*}
\FloatBarrier

\textbf{Exit for the ATP as a whole:} every ship lot (or defined sample) has
traceable pass data against versioned limits tied to the requirements slice.
\textbf{Decision:} ship, hold, or reopen DVT limits.

\section{ANOVA gauge R\&R sketch}
\label{sec:mfg-anova-gauge}

This is a procedure-level sketch, not a statistics course. Worksheets,
acceptance tables, and attribute studies live in the AIAG measurement-systems
analysis reference~\cite{aiagmsa}. Why measurement-system analysis precedes
yield interpretation: \Cref{sec:gauge-rr}.

A crossed variable study typically uses $n$ parts that span the process range,
$k$ appraisers (operators, stations, or shifts), and $r$ replicates of each
part$\times$appraiser cell. An ANOVA random-effects model partitions the total
observed variance into part-to-part, repeatability (equipment / method
variation under fixed appraisal conditions), reproducibility (appraiser or
condition contribution, often including part$\times$appraiser interaction), and
gauge R\&R as the combined measurement contribution.

Report the study with the metric, reference plane, units under test, and the
decision the measurement must support. Common summary ratios include gauge R\&R
as a fraction of total observed variation or of the tolerance width, and the
number of distinct categories the measurement system can resolve. Treat those
ratios as decision aids for a named ATP row, not as a universal green-light
threshold for every optical metric. Destructive or non-replicable tests need
alternate designs; do not force a crossed replicate study that the physics
cannot support.

\section{Gauge R\&R and station-correlation template}
\label{sec:mfg-gauge-template}

Use this checklist when teaching or auditing a production station. Background
and ANOVA sketch: \Cref{sec:gauge-rr,sec:mfg-anova-gauge}.

\begin{itemize}
\item Define the metric, reference plane, and units under test (good, marginal,
      failing).
\item Repeatability: same station, operator, and unit; report within-station
      spread.
\item Reproducibility: across stations, shifts, operators, and fixtures.
\item Bias: production station versus reference lab (DCA, BERT, OSA, power
      meter).
\item Stability: golden-unit trend over time; custody, recertification, and
      retirement criteria for the golden.
\item False reject / false accept: guardband from measurement uncertainty.
\item Station correlation offset and spread; decision rule for when a station
      is held.
\item CMIS or embedded monitor correlation to the same bench instruments when
      telemetry is used as a production proxy.
\item ANOVA gauge R\&R study design and reported variance components when a
      quantitative measurement-system claim is required (\Cref{sec:mfg-anova-gauge}).
\end{itemize}

\section{FAIR and supplier evidence checklist}
\label{sec:mfg-fair-checklist}

When tooling, epi, assembly site, driver/TIA silicon, TEC vendor, FAU epoxy, or
CMIS firmware changes, require a defined first-article package before open PO
volume:

\begin{itemize}
\item Change notice: what changed, why, affected lots, and effective date.
\item Pre- and post-change genealogy for representative units.
\item Measurement-system status: gauge R\&R or station correlation still valid
      for the affected ATP rows.
\item Multi-lot first-pass yield and parameter distributions for the changed
      population.
\item Capability and guardband check on critical parameters.
\item Qualification re-entry depth when life or environmental mechanisms are
      affected (\Cref{ch:reliability,app:reliability-reference}).
\item Incoming sample plan and SPC metrics with owners and reaction plans.
\item ECO and firmware version control matching the production reference freeze
      in \Cref{ch:manufacturing}.
\end{itemize}

\section{Sample control-plan mapping}
\label{sec:mfg-control-plan-map}

Requirements and the production control plan form the contract. Every
requirement needs a named evidence source and owner, but that evidence may be
ATP, a sampled audit, qualification, process control, supplier evidence, or
fleet monitoring.

\begin{enumerate}
\item \textbf{Requirements / PRD slice for the laser path:} fill
      \Cref{tab:laser-prd,sec:laser-reqs} (power class, grid, RIN@ORL, SMSR,
      derating, CMIS, FIT). Version it with the production control plan.
\item \textbf{Acceptance test plan (ATP):} the measurable tests that prove those
      requirements on every ship lot (or on a defined sample) when every-unit or
      sampled production test is the right control. Map each ATP line to a
      requirement; map life claims to qualification where appropriate
      (\Cref{sec:gr468,tab:atp-laser}).
\end{enumerate}

\heuristic{If a requirement has no evidence source, owner, or reaction plan, it
is a wish. ATP is only one possible control. If an ATP line has no requirement,
it is cost without a decision.}

\section{Rapid Interview Checks}

These are short prompts for self-test. Worked spoken answers live in
\Cref{sec:manufacturing-interview-qa}.

\paragraph{Prompt.} What must be frozen before a production-intent build?\\
\textit{Check.} Hardware/BOM, suppliers, firmware, calibration, recipes,
fixtures, test software, limits, and explicit deviations.

\paragraph{Prompt.} Why clear the measurement system before yield?\\
\textit{Check.} Otherwise station drift, golden aging, or bias masquerades as
product or supplier failure (\Cref{sec:gauge-rr}).

\paragraph{Prompt.} What proves an ATP screen detects defects?\\
\textit{Check.} Naturally failing units or controlled fault injection across
severity; detection probability and false rejects. Passing good units is not
enough.

\paragraph{Prompt.} When is a field fail a manufacturing escape?\\
\textit{Check.} Only when evidence ties the mechanism to a preventable
production or test-control gap. Otherwise wear-out, interop, install, service,
software, or residual latent risk (\Cref{sec:escaped-defects,ch:failure-modes}).
